% Generated draft: requires mathematical and source review.
\documentclass[11pt]{article}
\usepackage[margin=25mm]{geometry}
\usepackage{fontspec}
\setmainfont{Latin Modern Roman}
\usepackage{amsmath,amssymb,mathtools,unicode-math}
\setmathfont{Latin Modern Math}
\usepackage{enumitem,needspace,xurl,hyperref}
\setlength{\parindent}{0pt}
\setlength{\parskip}{5pt}
\setlength{\emergencystretch}{3em}
\newcommand{\N}{\mathbb{N}}
\newcommand{\Z}{\mathbb{Z}}
\newcommand{\Q}{\mathbb{Q}}
\newcommand{\R}{\mathbb{R}}
\newcommand{\C}{\mathbb{C}}
\newcommand{\F}{\mathbb{F}}
\newcommand{\A}{\mathbb{A}}
\newcommand{\PP}{\mathbb{P}}
\newcommand{\E}{\mathbb{E}}
\newcommand{\eps}{\varepsilon}
\newcommand{\dd}{\,\mathrm{d}}
\begin{document}
\section*{UnsolvedMath ID: 30000677}
\textbf{Source problem code:} catalogue entry
\subsection{Definitions and mathematical statement}
Let \(L\) be algebraically closed of characteristic \(p>0\), and identify \(\overline{\F}_p\) with the subfield of elements of \(L\) algebraic over \(\F_p\). Let \(G/L\) be semiabelian: a connected commutative algebraic group fitting into an exact sequence
\[
 0\longrightarrow T\longrightarrow G\longrightarrow A\longrightarrow0,
\]
where \(T\) is a torus and \(A\) an abelian variety. Let \(X\subseteq G\) be an irreducible subvariety and \(\Lambda\leq G(L)\) a finite-rank subgroup. Under the ordinary abelian-group rank convention this means \(\dim_\Q(\Lambda\otimes_\Z\Q)<\infty\), not finite generation; the source does not give a different rank convention or restrict denominators or torsion.

Assume \(X(L)\cap\Lambda\) is Zariski dense in \(X\). The full positive-characteristic Mordell--Lang question in the record asks whether there exist
\[
 g\in G(L),\quad G'\leq G\text{ an algebraic subgroup},\quad
 G_0/\overline{\F}_p\text{ semiabelian},\quad
 X_0\subseteq G_0\text{ defined over }\overline{\F}_p,
\]
and a surjective algebraic group homomorphism over \(L\),
\[
 h:G'\longrightarrow (G_0)_L,
\]
such that
\[
 X=g+h^{-1}((X_0)_L).
\]
Here the subscript \(L\) denotes base change, making the inverse image and translation in \(G\) precise. This is the source's definition of special. The translation \(g\), subgroup \(G'\) and map \(h\) are not required to descend to the constant field; the specified constant objects are \(G_0,X_0\).

The formula is retained as equality with the full inverse image, not weakened to containment or to a selected irreducible component. For a variety interpretation in characteristic \(p\), it is naturally an equality of reduced subvarieties/underlying closed sets; the source does not specify a nonreduced scheme-theoretic convention. No connectedness of \(G'\), separability of \(h\), or properness of \(G'\) is appended. The conclusion is not replaced by an ordinary subgroup coset, which would discard its constant-variety component.

The hypotheses range over the full stated finite-rank class, rather than only finitely generated groups or prime-to-\(p\) divisions. The literature triage checked 2026-08-21 reports no authoritative complete resolution of this extracted claim located in its search. Its open label and reported official-source verification are retained as dated metadata, not independently checked current status.
\subsection{Short English statement}
Does a Zariski-dense intersection with a finite-rank subgroup force the stated constant-variety pullback description of \(X\)? The characteristic-\(p\) constant data and full inverse-image equality are preserved. Finite rank is not silently restricted to finite generation, and the map and translation need not be constant-defined.
\subsection{Sources}
\begin{itemize}
\item Local UnsolvedMath registry, permanent ID 30000677, source code catalogue entry, ``Mordell--Lang for Semiabelian Varieties in Positive Characteristic,'' including the original full-conjecture designation, explicit special-variety definition and verification metadata.
\item \emph{Model Theory and Groups (2007)}, Oberwolfach Reports, pp.\ 83--138, DOI 10.4171/owr/2007/02: \url{https://doi.org/10.4171/owr/2007/02}. This is also the sole supplied literature reference.
\item Literature triage checked 2026-08-21; record updated 2026-09-16. No external source has been visited or independently checked here.
\end{itemize}
\end{document}
